\DocumentMetadata{
  pdfversion=2.0,
  pdfstandard=ua-2,
  lang=en-US,
  tagging=on,
  tagging-setup={math/setup=mathml-SE,tabsorder=structure}
}
\documentclass[11pt,letterpaper]{article}
\usepackage[margin=0.68in,includefoot]{geometry}
\usepackage{newcomputermodern}
\usepackage{amsmath,amscd,mathtools}
\usepackage{tikz,adjustbox}
\providecommand{\square}{\mdlgwhtsquare}
\usepackage[hidelinks]{hyperref}
\hypersetup{
  pdftitle={Algebra PhD Exam, May 2010},
  pdfauthor={University of Florida Department of Mathematics},
  pdfsubject={Graduate mathematics examination},
  pdfdisplaydoctitle=true,
  bookmarksopen=true
}
\setcounter{secnumdepth}{0}
\setlength{\parindent}{0pt}
\setlength{\parskip}{0.28em}
\setlength{\emergencystretch}{3em}
\setlength{\leftmargini}{2.15em}
\setlength{\leftmarginii}{2.65em}
\setlength{\leftmarginiii}{3em}
\pagestyle{empty}
\raggedbottom
\allowdisplaybreaks
\NewTaggingSocketPlug{block/list/label}{examproblem}{%
  \tagstructbegin{tag=Lbl}\tagstructend%
  \tagstructbegin{tag=\LBody}%
  \tagstructbegin{tag=\ProblemHeadingTag,title-o={\ProblemHeadingText}}%
  \tagmcbegin{tag=\ProblemHeadingTag,actualtext={\ProblemHeadingText}}%
  #2%
  \tagmcend%
  \tagstructend%
}
\newenvironment{examproblems}[2]{%
  \def\ProblemHeadingTag{#2}%
  \AssignTaggingSocketPlug{block/list/label}{examproblem}%
  \begin{enumerate}%
  \setcounter{enumi}{#1}%
  \renewcommand{\labelenumi}{\textbf{\theenumi.}}%
  \setlength{\topsep}{0.35em}%
  \setlength{\partopsep}{0pt}%
  \setlength{\itemsep}{0.48em}%
  \setlength{\parsep}{0.18em}%
}{\end{enumerate}}
\newenvironment{examsubparts}{%
  \AssignTaggingSocketPlug{block/list/label}{default}%
  \begin{enumerate}%
  \setlength{\topsep}{0.18em}%
  \setlength{\partopsep}{0pt}%
  \setlength{\itemsep}{0.16em}%
  \setlength{\parsep}{0.1em}%
}{\end{enumerate}}
\newenvironment{examinstructions}{%
  \par\begingroup\itshape
}{%
  \par\endgroup\vspace{0.18em}
}
\makeatletter
\renewcommand\subsection{\@startsection{subsection}{2}{0pt}%
  {-1.25ex plus -.25ex minus -.1ex}%
  {0.55ex plus .1ex}%
  {\normalfont\large\bfseries}}
\renewcommand{\@maketitle}{%
  \newpage\null\vskip -1em%
  {\footnotesize\bfseries
    \href{https://www.ufl.edu/}{University of Florida}, \href{https://math.ufl.edu/}{Department of Mathematics}\par}%
  \vskip -0.5em%
  \tagstructbegin{tag=H1,title={Algebra PhD Exam, May 2010}}%
  \tagpdfparaOff%
  \tagmcbegin{tag=H1}%
    {\LARGE\bfseries\href{https://gma.math.ufl.edu/exams/algebra-phd/}{Algebra PhD Exam}, May 2010\par}%
  \tagmcend%
  \tagpdfparaOn%
  \tagstructend%
  \vskip 0.25em%
  \tagmcbegin{artifact}%
    \hrule height 1pt%
  \tagmcend%
  \vskip 1em%
}
\makeatother
\title{Algebra PhD Exam, May 2010}
\author{}
\date{}
\begin{document}
\maketitle
\thispagestyle{empty}
\begin{examinstructions}
Answer seven problems. Write your answers clearly in complete English sentences. You may quote results (within reason) as long as you state them clearly.
\end{examinstructions}
\begin{examproblems}{0}{H2}
\def\ProblemHeadingText{Problem 1.}
\item
Prove that every field has an algebraic closure.
\def\ProblemHeadingText{Problem 2.}
\item
This problem has three parts.
\begin{examsubparts}
\item[\textnormal{(a)}]
Determine a splitting field~\(K\) for the polynomial \(X^8-1\in\mathbb{Q}[X]\).
\item[\textnormal{(b)}]
Determine \(\operatorname{Gal}(K/\mathbb{Q})\) and give an explicit description of the action of \(\operatorname{Gal}(K/\mathbb{Q})\) on~\(K\).
\item[\textnormal{(c)}]
For each subgroup~\(H\) of \(\operatorname{Gal}(K/\mathbb{Q})\) find the subfield of~\(K\) fixed by~\(H\).
\end{examsubparts}
\def\ProblemHeadingText{Problem 3.}
\item
Let~\(E\) and~\(F\) be fields such that \(\operatorname{char}(E)\ne\operatorname{char}(F)\). Let~\(V\) be a vector space over~\(E\) and let~\(W\) be a vector space over~\(F\). Prove that \(V\otimes_{\mathbb{Z}}W\) is trivial.
\def\ProblemHeadingText{Problem 4.}
\item
This problem has four parts.
\begin{examsubparts}
\item[\textnormal{(a)}]
Give the definition of a covariant functor from a category~\(C\) to a category~\(D\).
\item[\textnormal{(b)}]
Give the definition of a concrete category.
\item[\textnormal{(c)}]
Define what it means for an object~\(F\) in a concrete category to be free on a set~\(S\).
\item[\textnormal{(d)}]
Given a set~\(S\), construct a free object on~\(S\) in the category of pointed sets.
\end{examsubparts}
\def\ProblemHeadingText{Problem 5.}
\item
This problem has two parts.
\begin{examsubparts}
\item[\textnormal{(a)}]
Define what it means for an~\(R\)-module~\(M\) to be injective.
\item[\textnormal{(b)}]
Let \(\{M_\lambda:\lambda\in\Lambda\}\) be a nonempty collection of injective~\(R\)-modules. Prove that the product \(\prod_{\lambda\in\Lambda}M_\lambda\) is injective.
\end{examsubparts}
\def\ProblemHeadingText{Problem 6.}
\item
Let~\(F\) be a free group on a set~\(S\) with \(|S|>1\). Prove that~\(F\) is not solvable.
\def\ProblemHeadingText{Problem 7.}
\item
Let~\(R\) be a commutative Noetherian ring with~\(1\) and let \(I\subset R\) be an ideal. Prove that~\(I\) has a primary decomposition.
\def\ProblemHeadingText{Problem 8.}
\item
Let~\(S\) be a commutative ring with~\(1\), let~\(R\) be a subring of~\(S\) which contains \(1_S\), and let \(x\in S\). Suppose there is a subring \(T\subset S\) such that \(T\supset R[x]\) and~\(T\) is finitely generated as an~\(R\)-module. Prove that~\(x\) is integral over~\(R\).
\def\ProblemHeadingText{Problem 9.}
\item
Prove Noether’s Normalization Lemma: Let~\(k\) be a field and let~\(A\) be a finitely generated~\(k\)-algebra with~\(1\). Then there are \(y_1,\ldots,y_q\in A\) such that \(y_1,\ldots,y_q\) are algebraically independent over~\(k\) and~\(A\) is integral over \(k[y_1,\ldots,y_q]\).
\def\ProblemHeadingText{Problem 10.}
\item
This problem has two parts.
\begin{examsubparts}
\item[\textnormal{(a)}]
Prove that every nonzero proper ideal in a Dedekind domain is a product of prime ideals.
\item[\textnormal{(b)}]
Give an example of an integral domain~\(R\) and a nonzero proper ideal \(I\subset R\) such that~\(I\) cannot be expressed as a product of prime ideals.
\end{examsubparts}
\def\ProblemHeadingText{Problem 11.}
\item
Let~\(G\) be a finite group of order~\(n\) and let~\(F\) be a field.
\begin{examsubparts}
\item[\textnormal{(a)}]
Prove that if \(\operatorname{char}(F)\nmid n\) then for every \(F[G]\)-module~\(V\) and every \(F[G]\)-submodule \(W\subset V\) there is an \(F[G]\)-submodule \(W'\subset V\) such that \(V=W\oplus W'\).
\item[\textnormal{(b)}]
Show that if \(\operatorname{char}(F)\mid n\) then there exists an \(F[G]\)-module~\(V\) and an \(F[G]\)-submodule \(W\subset V\) such that there is no \(F[G]\)-submodule \(W'\subset V\) such that \(V=W\oplus W'\).
\end{examsubparts}
\end{examproblems}
\end{document}
